% Author: Adam | Module: AI Job Interview Simulator | Project: SmartCareer AI
\section{Ingénierie des prompts}
\label{sec:prompts}

C'est le cœur du module : aucun fine-tuning n'étant conduit ici (le Module 3
utilise le modèle tel que quantifié), toute la qualité de sortie provient des
prompts et du pipeline qui les entoure. Quatre prompts sont produits par du
code : trois prompts métier (\texttt{adam\_questions}, \texttt{adam\_evaluate},
\texttt{adam\_improve}) et un prompt de réparation (\texttt{repair}).

\subsection{Principes de conception}

\begin{enumerate}[label=\textbf{P\arabic*.},leftmargin=3.4em,itemsep=3pt]
  \item \textbf{Une seule sortie : du JSON.} Chaque prompt se termine par un
        contrat de format explicite (\texttt{FORMAT : renvoie exactement \{...\}})
        et par le rappel \texttt{FORMAT\_REMINDER}. La grammaire
        \texttt{format=} d'Ollama est une aide, jamais une garantie
        (§\ref{sec:modele-format}).
  \item \textbf{Barème à ancres.} Un 7\,Md sans échelle de référence note tout
        entre 7 et 8 : le prompt d'évaluation définit cinq paliers de 0 à 10
        avec, pour chacun, ce qui le caractérise, et quatre critères
        indépendants.
  \item \textbf{Les données entre chevrons ne sont pas des instructions.} Tout
        contenu externe (offre, CV, réponse du candidat, bloc RAG) est enveloppé
        dans \verb|<label>…</label>| et une clause interdite au modèle d'en
        exécuter le contenu (P3 de §\ref{sec:sec-inject}).
  \item \textbf{Interdiction d'inventer, avec alternative matérialisée.} Là où
        l'information manque, le module demande un marqueur
        \verb|[à compléter : …]| plutôt qu'un remplissage : une réponse de
        30~mots balisée est un succès, une histoire inventée de 120~mots est une
        faute.
  \item \textbf{Langue de sortie découplée de la langue des ressources.} La base
        documentaire RAG est en anglais ; une clause impose la langue demandée
        « même si un document de référence fourni est rédigé dans une autre
        langue » (§\ref{sec:prompts-problemes}, problème~2).
  \item \textbf{L'exemple ne fournit que la forme.} Les blocs few-shot sont
        marqués comme squelettes et une liste de leurs mots-signatures est
        explicitement interdite en sortie.
  \item \textbf{Budget de prompt borné, rognage ordonné.} Les sections portent
        un poids ; au-delà du budget, les plus faibles sont rognées en premier
        (\texttt{base\_adam.fit}, §\ref{sec:prompts-assembleur}).
\end{enumerate}

\subsection{Un assembleur commun}
\label{sec:prompts-assembleur}

\file{prompts/\brk{}base\_adam.py} centralise quatre constantes et quatre fonctions.
Les constantes, reproduites verbatim (elles constituent la partie commune des
trois \texttt{SYSTEM}) :

\medskip
\noindent\textbf{\texttt{SYSTEM\_HEADER}} (674~car.) --- rôle, contrat JSON,
interdiction d'inventer, consigne de langue.
\input{\adamroot sections/lst-base-system_header}

\smallskip
\noindent\textbf{\texttt{IGNORE\_INJECTED}} (238~car.) --- clause anti-injection
apposée à tout \texttt{SYSTEM} par \texttt{compose()}.
\input{\adamroot sections/lst-base-ignore_injected}

\smallskip
\noindent\textbf{\texttt{JSON\_GUARDRAIL}} (198~car.) --- n'est utilisé que dans
le prompt de réparation, via \texttt{schema\_contract()}.
\input{\adamroot sections/lst-base-json_guardrail}

\smallskip
\noindent\textbf{\texttt{FORMAT\_REMINDER}} (67~car.) --- rappel final du prompt
normal.
\input{\adamroot sections/lst-base-format_reminder}

Trois mécaniques complètent l'assemblage.

\begin{itemize}[itemsep=2pt]
  \item \texttt{block(label, text, limit)} : \texttt{clip} (borne en
        caractères) $\to$ \texttt{scrub\_text} (suppression e-mail, téléphone,
        URL) $\to$ \texttt{wrap\_content} (neutralisation d'un gabarit
        étranger) $\to$ enveloppe \verb|<label>|. Les prompts sont donc sûrs
        \emph{par construction} ; le contrôle du client
        (\texttt{assert\_anonymized}) n'est qu'un garde-fou contre une
        régression (§\ref{sec:sec-pii}).
  \item \texttt{fit(sections, budget)} : rogne les sections de plus faible
        poids jusqu'à tenir dans le budget. Budgets effectifs, calculés depuis
        \file{config\_adam.py} : 8500~car. pour les questions, 8900 pour
        l'évaluation, 14\,100 pour l'amélioration.
  \item \texttt{compose(system, user, schema)} : assemble
        \verb|[INST] <<SYS>> … <</SYS>> …  [/INST]| \emph{sans} recopier le
        schéma en texte (il voyage dans \texttt{format=}) --- décision prise
        pour la taille, cf. tableau~\ref{tab:taille-prompts}.
\end{itemize}

Une lecture attentive des listings de cette section demande deux conventions :
les fragments sont extraits du \emph{code source}, donc les reprises de ligne
avec guillemets adjacents sont la concaténation implicite de chaînes Python, et
\verb|{{ }}| est l'échappement d'une f-string --- le modèle reçoit
\verb|{ }|.

\subsection{Prompt n°1 --- génération des questions}
\label{sec:prompt-questions}

\textbf{But.} Produire exactement $n$ questions pertinentes pour \emph{ce}
poste, chacune accompagnée de son barème privé (\texttt{expected\_points},
2 à 4 éléments attendus) qui servira à l'évaluation.

\textbf{Entrées.} \texttt{job\_title}, \texttt{seniority},
\texttt{question\_type}, \texttt{n\_questions}, \texttt{language} ; optionnels
\texttt{company}, \texttt{job\_description}, \texttt{cv\_digest\_text},
\texttt{gaps} (écarts offre/profil calculés sans LLM), \texttt{rag\_context}.

\textbf{Ordre des sections} (poids décroissant, §\ref{sec:prompts-assembleur}) :
TÂCHE avec la répartition catégorielle calculée \textbf{(100)} --- poste visé
\textbf{(80)} --- profil anonymisé ou mention « aucun CV » \textbf{(30)} ---
manques \textbf{(20)} / compétences déjà visibles \textbf{(10)} --- références
RAG \textbf{(15)} --- consigne de langue \textbf{(90)} --- exemple
\textbf{(60)}.

Le calcul de la répartition mérite une ligne : \texttt{min\_majority =
(2n + 2) // 3}, soit $\lceil 2n/3 \rceil$. Pour $n=3$, le prompt exige
« 2 questions au minimum sur 3 en catégorie technical ». Sans cette arithmétique
le modèle alternait librement et une demande « technique » produisait
fréquemment trois questions comportementales.

\textbf{\texttt{SYSTEM} du prompt de questions, verbatim :}
\input{\adamroot sections/lst-questions-system}

\textbf{Bloc few-shot, verbatim :}
\input{\adamroot sections/lst-questions-example}

\textbf{Schéma JSON attendu} (\texttt{QUESTION\_SET\_Schema}, transmis en
\texttt{format=} ; un élément \texttt{questions} porte \texttt{question},
\texttt{category} dans \{\texttt{technical}, \texttt{project}, \texttt{hr}\},
\texttt{focus}, \texttt{expected\_points}) :
\input{\adamroot sections/lst-schema-question-set}

\textbf{Justification des choix.} « une question = une seule chose, 40 mots
max, à la 2\textsuperscript{e} personne » corrige des questions fleuves de trois
sous-questions. \texttt{expected\_points} est demandé ici plutôt qu'à
l'évaluation pour que le barème soit fixé \emph{avant} de voir la réponse ---
sinon le modèle ajuste le barème à la réponse et la note devient circulaire.
L'exemple est en français même pour une séance anglaise : c'est un choix
assumé (P5), la consigne de langue se chargeant de la traduction, ce qui évite
d'exemple en double et ~780~caractères de prompt.

\subsection{Prompt n°2 --- évaluation d'une réponse}
\label{sec:prompt-evaluate}

\textbf{But.} Noter une réponse sur 10 avec quatre critères séparés, un
verdict formatif, et un diagnostic de \emph{type} d'échec.

\textbf{Entrées.} \texttt{question}, \texttt{category}, \texttt{expected\_points}
(repris du prompt 1, tronqué à~6), \texttt{answer} (borné à
\texttt{MAX\_ANSWER\_CHARS}=4000), \texttt{language}, \texttt{job\_title},
\texttt{seniority}, \texttt{rubric\_context} (RAG).

Deux décisions structurelles : la question et sa catégorie sont rappelées dans
le prompt (le module n'a pas de mémoire de conversation, §\ref{sec:modele-raw}) ;
le barème du recruteur est injecté comme « éléments attendus », ce qui abaisse
fortement la variance des notes.

\textbf{\texttt{SYSTEM} du prompt d'évaluation, verbatim :}
\input{\adamroot sections/lst-evaluate-system}

\textbf{Bloc few-shot, verbatim :}
\input{\adamroot sections/lst-evaluate-example}

\textbf{Schéma JSON attendu} (score 0--10 décimal, quatre critères indépendants,
forces et faiblesses en tableaux de chaînes, verdict libre) :
\input{\adamroot sections/lst-schema-evaluation}

\textbf{Justification.} Le barème à cinq paliers et les quatre critères
indépendants viennent d'un constat : sans ancres, le modèle concentre ses notes
entre 7 et 8 (quantifié en §\ref{sec:tests-calibration}, où l'écart entre une
mauvaise et une bonne réponse passe à 4,0 points après réglage).
L'interdiction « juge uniquement ce qui est écrit » répond au même problème sous
un autre angle : sans elle le modèle créditait des compétences absentes de la
réponse parce que le CV les mentionnait.

\subsection{Prompt n°3 --- réponse améliorée}
\label{sec:prompt-improve}

\textbf{But.} Fournir un \emph{modèle} de réponse que le candidat réécrit avec
ses vrais faits : structure STAR pour le comportemental, démonstration de
raisonnement pour le technique, plus 2--3 conseils pour l'oral.

\textbf{Entrées.} \texttt{question}, \texttt{category}, \texttt{answer},
\texttt{language}, \texttt{evaluation} (score, critères faibles, jusqu'à 4
lacunes), \texttt{cv\_digest\_text}, \texttt{job\_title}, \texttt{seniority},
\texttt{technique\_notes} (RAG). Les lacunes de l'évaluation sont injectées en
« points à corriger en priorité » : l'amélioration est appelée sur demande, après
évaluation, et s'appuie sur son diagnostic.

\textbf{\texttt{SYSTEM} du prompt d'amélioration, verbatim :}
\input{\adamroot sections/lst-improve-system}

\textbf{Bloc few-shot, verbatim :}
\input{\adamroot sections/lst-improve-example}

\textbf{Schéma JSON attendu} :
\input{\adamroot sections/lst-schema-improvement}

\textbf{Justification.} La clause « reste sur le SUJET et la MÊME situation » a
été ajoutée après observation d'une dérive narrative (question sur la
conteneurisation réécrite en anecdote de débogage de logs). L'encadrement de
longueur (80--120 mots \emph{seulement si} la matière existe) tient à la
latence : en CPU, ce prompt produit environ 4 tokens/s et une réponse de
200~mots double le temps d'attente pour aucun gain pédagogique (§
\ref{sec:perf-latence}).

\subsection{Prompt de réparation}
\label{sec:prompt-repair}

Déclenché quand la première sortie est rejetée (parse JSON ou Pydantic). Il
n'y a \textbf{qu'une} tentative de réparation, puis repli
(§\ref{sec:prompts-validation}).

Sa particularité : c'est le seul endroit où le schéma JSON est recopié en
toutes lettres, parce que le but est précisément de montrer la forme attendue.
\texttt{build()} reçoit le prompt d'origine, la sortie fautive (bornée à
1600~car.) et le motif de rejet (borné à 400~car.). Comme \texttt{raw=true}
ignore la mémoire de contexte, la réparation est présentée comme un second tour
de conversation : le \texttt{[/INST]} d'origine est refermé, puis une nouvelle
requête d'ouverture est ajoutée --- le texte d'en-tête est reproduit verbatim :
\input{\adamroot sections/lst-repair-repair_header}

Sur l'exemple de calibration, le prompt de réparation ainsi obtenu mesure
5539~caractères. Il n'a \textbf{jamais été exécuté en conditions réelles} : les
10 appels de la calibration et les 14 appels du parcours Playwright ont tous
été valides dès le premier essai (§\ref{sec:tests-calibration}) --- le chemin
est couvert par 21 tests unitaires à client simulé, pas par la production.

\subsection{Problèmes rencontrés et corrections}
\label{sec:prompts-problemes}

Six défauts ont été observés, corrigés, puis revérifiés. Pour chacun, la
columnne « preuve » renvoie au texte présent dans le code.

\begin{small}
\begin{longtable}{@{}>{\cellflush}p{2.5cm}
                  >{\cellflush}p{4.3cm}
                  >{\cellflush}p{4.4cm}
                  >{\cellflush}p{3.4cm}@{}}
\toprule
\textbf{Défaut} & \textbf{Observé} & \textbf{Correction} & \textbf{Preuve dans le code} \\
\midrule
\endhead
\bottomrule
\endlastfoot
Défaut 1 --- fuite de l'exemple & La réponse améliorée d'une question Docker parlait d'« inscriptions », de « pool de connexions » et de « taux d'échec » : des éléments de l'exemple few-shot. & L'exemple est devenu un squelette marqué « EXEMPLE DE FORMAT UNIQUEMENT --- aucun contenu réel à recopier », et une règle INTERDIT nommément ses mots-signatures sauf s'ils figurent dans la réponse du candidat. & \file{adam\_improve.py}:33--35 et \texttt{EXAMPLE}:53 \\
\addlinespace
Défaut 2 --- anglais dans une sortie française & Sur une séance en français, les \texttt{strengths}/\texttt{weaknesses} arrivaient en anglais dès qu'un bloc RAG était injecté : la base documentaire est anglaise et le modèle alignait sa langue dessus. & Clause dans \texttt{SYSTEM\_HEADER} : « Rédige TOUTES les valeurs de texte \ldots{} dans la langue demandée, même si un document de référence fourni est rédigé dans une autre langue », complétée par \texttt{language\_instruction()} et, côté évaluation, « ne recopie aucun document de référence ». & \file{base\_adam.py}:36--38, 74--82 \\
\addlinespace
Défaut 3 --- résultats inventés & Une amélioration concluait par « 35\,\% de temps de réponse en moins » alors que la réponse du candidat ne contenait aucun chiffre. & Interdiction formelle d'ajouter un fait (entreprise, techno, chiffre, durée, résultat) absent de la réponse ou du profil, et interdiction de conclure sur un résultat inventé ; le manquement devient un marqueur \verb|[à compléter : …]|. & \file{adam\_improve.py}:27--32 \\
\addlinespace
Défaut 4 --- hors-sujet confondu avec lacune & Une réponse juste mais à côté se voyait reprocher un manque de connaissances techniques ; le conseil qui en découlait (« révisez le SQL ») était hors cible. & Le \texttt{SYSTEM} d'évaluation distingue nommément HORS SUJET (\texttt{relevance} basse, le reste peut être juste) et MANQUE DE CONNAISSANCES (\texttt{relevance} correcte, \texttt{technical\_accuracy} basse), avec interdiction explicite de les confondre ; le \texttt{verdict} doit dire lequel. & \file{adam\_evaluate.py}:37--42 \\
\addlinespace
Défaut 5 --- prompt trop volumineux & Le prompt de questions dépassait 8500~caractères : préfill lent en CPU et fenêtre de 8\,192~jetons occupée aux deux tiers. & Suppression du schéma en texte dans les prompts normaux (il reste dans \texttt{format=}), exemples réduits, consignes condensées, et rognage pondéré par \texttt{fit()}. Voir tableau~\ref{tab:taille-prompts}. & \file{base\_adam.py}:76--86, \texttt{compose()} \\
\addlinespace
Défaut 6 --- fuite d'identité par le titre du CV & Le titre d'une génération du Module 1 est un nom de fichier librement saisi ; il peut contenir le nom du candidat. Transmis au prompt, il exposait une identité dans le journal et dans les \texttt{rag\_sources}. & Le titre n'est plus jamais exposé au Module 3 : \texttt{get\_cv\_bundle()} ne le renvoie pas ; seule la fonction recherchée du profil (\texttt{target\_title}) est conservée. & \file{cv\_bridge.py}:57--61 (commentaire explicite) \\
\end{longtable}
\end{small}

\subsection{Taille des prompts avant et après l'optimisation}
\label{sec:prompts-taille}

\begin{small}
\begin{table}[H]
\centering
\caption{Caractères de prompt, avant la passe d'optimisation puis après
(réduction entre crochets). Les tokens de prompt sont relevés dans les journaux
du parcours live, sur les entrées réelles de la démonstration.}
\label{tab:taille-prompts}
\begin{tabular}{@{}l r r r r r@{}}
\toprule
\textbf{Étape} & \textbf{Avant} & \textbf{Après} & \textbf{Réduc.} & \textbf{Tok. prompt (live)} & \textbf{\texttt{num\_predict}} \\
\midrule
Questions   & 8\,520 & 6\,078 & $-28,7\,\%$ & 1\,797 & 630 ($n=3$) \\
Évaluation  & 7\,478 & 5\,781 & $-22,7\,\%$ & 1\,792 & 420 \\
Amélioration & 10\,524 & 6\,772 & $-35,7\,\%$ & 2\,029 & 520 \\
\bottomrule
\end{tabular}
\end{table}
\end{small}

Trois chiffres à ne pas confondre. (a)~Le tableau ci-dessus compare l'état
\emph{initial} et l'état \emph{après} la passe d'optimisation, sur le contexte
d'une séance réelle. (b)~Un même jeu de prompts, re-mesuré sur un jeu
d'entrée de contrôle unique (offre Java/Angular de 213~car., CV synthétique de
213~car., deux fragments RAG de 263~car.), donne aujourd'hui : questions
4427~car. (système 1975, utilisateur 2452, 636~mots), évaluation 4812~car.
(2957 / 1855, 694~mots), amélioration 5521~car. (3097 / 2424, 835~mots),
réparation 5539~car. La différence avec (a) vient uniquement de la taille des
entrées, l'entrée candidat étant bornée à 4000~car. (c)~Le budget de rognage
(8500 / 8900 / 14\,100) n'est pas une taille atteinte : c'est le plafond au-delà
duquel \texttt{fit()} coupe les sections de plus faible poids.

\subsection{Pipeline de validation de la sortie}
\label{sec:prompts-validation}

\begin{enumerate}[label=\arabic*.,leftmargin=*,itemsep=2pt]
  \item \textbf{Nettoyage de surface} (\texttt{\_clean\_text}) : retrait d'un
        éventuel bloc \verb|<think>| et des fences Markdown.
  \item \textbf{Analyse JSON tolérante} (\texttt{\_parse\_json}) : \texttt{json.loads}
        direct ; en cas d'échec, nouvelle tentative sur la tranche entre le premier
        \texttt{\{} et le dernier \texttt{\}} --- ce qui sauve une sortie encadrée
        de texte. Un objet non dict est rejeté.
  \item \textbf{Validation Pydantic} : \texttt{model\_validate\_json} sur le
        schéma de l'étape. Les rejets typiques : score hors de \texttt{ge=0, le=10},
        critère manquant, chaîne vide là où le schéma exige \texttt{min\_length}.
  \item \textbf{Réparation} : \textbf{une seule} tentative, température 0,1, avec
        le motif de rejet et le schéma en toutes lettres (§\ref{sec:prompt-repair}).
  \item \textbf{Repli déterministe} : \texttt{fallbacks.py}. Questions --- pool de
        questions génériques FR/EN par type ; Évaluation --- heuristique sur la
        forme (longueur, première personne, connecteurs, absence de repère chiffré),
        avec \texttt{score = None} ; Amélioration --- reformulation minimale.
  \item \textbf{Normalisation} : \texttt{\_normalize\_questions} indexe, tronque
        à la demande, complète à défaut les questions manquantes en les marquant
        \texttt{degraded: true}. \textbf{Jamais plus} que le nombre demandé.
\end{enumerate}

\textbf{Ce que « degraded » veut dire.} Le champ \texttt{\_meta.degraded} est un
\texttt{True}/\texttt{False} posé par \texttt{\_run}, et \texttt{reason} porte le
motif de rejet (tronqué à 200~car.). \texttt{degraded = true} signifie que
\textbf{la réponse affichée ne vient pas du modèle} : elle vient du repli. Le
degré de dégradation est visible côté interface (bandeau et marqueur sur la
carte) et le \texttt{score} vaut \texttt{None} dans le cas de l'heuristique ---
un choix volontaire : une note inventée par une heuristique serait une
contrefaçon. Une transport failure (Ollama injoignable, timeout, 429, 5xx) ne
dégrade \textbf{pas} : elle remonte en erreur HTTP (§\ref{sec:sec-erreurs}) pour
que l'utilisateur puisse réessayer, au lieu d'afficher un repli en silence.

\subsection{Les trois contrats de sortie}
\label{sec:prompts-schemas}

Les schémas reproduits plus haut (\texttt{QUESTION\_SET\_SCHEMA},
\texttt{EVALUATION\_SCHEMA}, \texttt{IMPROVEMENT\_SCHEMA}, fichiers
\texttt{lst-schema-*.tex}) sont écrits à la main dans
\file{schemas\_adam.py} et non générés depuis les modèles Pydantic, pour trois
raisons consignées dans l'en-tête de ce fichier.

\begin{enumerate}[label=\textbf{R\arabic*.},leftmargin=3.6em,itemsep=2pt]
  \item \textbf{Aucun \texttt{\$ref} ni \texttt{\$defs}.} La sortie de
        \texttt{model\_json\_schema()} de Pydantic v2 factorise les sous-objets par
        référence ; le schéma est aplati ici parce que la contrainte de décodage
        d'Ollama doit consommer un document autoportant.
  \item \textbf{Ni \texttt{enum}, ni \texttt{minimum}/\texttt{maximum}.} La
        grammaire de sortie ne les garantit pas et Mistral n'est pas listé parmi
        les modèles à sortie structurée ; les imposer risquait un blocage de
        décodage plutôt qu'une sortie valide. Les bornes sont donc
        \emph{re-calées côté Python} : \texttt{\_coerce\_score} accepte un nombre,
        une chaîne contenant un nombre (avec virgule) et \textbf{bornne à 0--10},
        avec arrondi au dixième ; une valeur non récupérable ou booléenne donne
        \texttt{None}. \texttt{ge=0, le=10} des modèles Pydantic reste la
        validation finale.
  \item \textbf{Forme courte.} Un schéma écrit à la main tient en 55~lignes là où
        la sérialisation Pydantic ajoute \texttt{title}, \texttt{description} et
        métadonnées --- du texte payé à chaque appel, alors que le même contrat est
        déjà énoncé en clair dans le \texttt{SYSTEM}.
\end{enumerate}

Le contrat n'est donc \emph{pas} une ceinture et des bretelles redondantes : le
\texttt{format} borne la \textbf{forme}, les validateurs Pydantic bornent la
\textbf{valeur}, et le repli bornent la \textbf{disponibilité}. Ce découpage est
ce qui permet d'afficher un score même quand le modèle n'a pas respecté
l'échelle demandée --- sans jamais inventer de note (cas \texttt{score = None}).
